% !TEX root = ../../../main.tex
% Sources: LEC03, IMG_9135--IMG_9138; Newton argument from LEC02, IMG_9118--IMG_9119.
% Keep the section heading with its opening discussion after the wide figure.
\clearpage
\section{Roots and rational powers}
\label{sec:analysis-i:roots-and-powers}

\newthought{An inverse function needs a bijection.} In the lecture's
notation, $\R_+=[0,\infty)$ includes zero. For $n\in\N$, we will prove
that $x\mapsto x^n$ is a bijection from $\R_+$ to itself. Only after that
proof can we define $x^{1/n}$ as its inverse.

\subsection{A difference of powers}
The finite geometric identity
\[
  1+r+\cdots+r^{n-1}=\frac{1-r^n}{1-r}\qquad(r\ne1)
\]
follows by multiplying the sum by $1-r$ and cancelling consecutive terms.
For $0\leq a<b$, substitute $r=a/b$ and multiply by $(b-a)b^{n-1}$.
This gives
\begin{equation}
  b^n-a^n=(b-a)\sum_{j=0}^{n-1}b^{n-1-j}a^j.
  \label{eq:analysis-i:power-difference}
\end{equation}
The sum is positive and has $n$ terms, each at most $b^{n-1}$. Consequently,
\begin{equation}
  0<b^n-a^n\leq n(b-a)b^{n-1}.
  \label{eq:analysis-i:power-bound}
\end{equation}
In particular, the power map is strictly increasing on $[0,\infty)$,
and hence injective. The upper bound controls how much the power changes
when the input moves a short distance.

\subsection{Completeness supplies every root}
\begin{theorem}
  \label{thm:analysis-i:positive-roots}
  For every $n\in\N$ and every $y\geq0$, there is a unique $x\geq0$
  such that $x^n=y$.
\end{theorem}
\begin{marginfigure}
  \centering
  % Equal units preserve the square's geometry at its printed size.
\begin{tikzpicture}[analysis diagram]
  \fill[black!10] (0,0) rectangle (2.2,2.2);
  \draw (0,0) rectangle (2.2,2.2);
  \node[analysis label] at (1.1,1.1) {$A=s^2$};
  \node[analysis label,below=6pt] at (1.1,0) {$s$};
  \node[analysis label,left=6pt] at (0,1.1) {$s$};
  \node[analysis label,below left=5pt] at (0,0) {$(0,0)$};
  \node[analysis label,above right=5pt] at (2.2,2.2) {$(s,s)$};
\end{tikzpicture}

  \caption[A square of prescribed area.]{The square $[0,s]\times[0,s]$ in
    $\R^2$ has area $A=s^2$. Every $A\geq0$ determines a unique
    nonnegative side length.}
  \label{fig:analysis-i:root-square}
\end{marginfigure}
\begin{marginfigure}
  \centering
  % An oblique view of the solid cube; hidden edges are omitted for clarity.
\begin{tikzpicture}[analysis diagram]
  \fill[black!6] (0,1.9) -- (0.8,2.5) -- (2.7,2.5) -- (1.9,1.9) -- cycle;
  \fill[black!16] (1.9,0) -- (2.7,0.6) -- (2.7,2.5) -- (1.9,1.9) -- cycle;
  \fill[black!10] (0,0) rectangle (1.9,1.9);
  \draw (0,0) rectangle (1.9,1.9);
  \draw (0,1.9) -- (0.8,2.5) -- (2.7,2.5)
    -- (2.7,0.6) -- (1.9,0);
  \draw (1.9,1.9) -- (2.7,2.5);
  \node[analysis label] at (0.95,0.95) {$V=s^3$};
  \node[analysis label,below=6pt] at (0.95,0) {$s$};
  \node[analysis label,left=6pt] at (0,0.95) {$s$};
  \node[analysis label,below right=6pt] at (2.3,0.3) {$s$};
\end{tikzpicture}

  \caption[A cube of prescribed volume.]{The cube
    $[0,s]\times[0,s]\times[0,s]$ in $\R^3$ has volume $V=s^3$.
    Every $V\geq0$ determines a unique nonnegative side length.}
  \label{fig:analysis-i:root-cube}
\end{marginfigure}
\begin{proof}
  Uniqueness follows from strict monotonicity. If $y=0$, take $x=0$.
  Suppose henceforth that $y>0$, and put
  \[
    S=\{t\geq0\mid t^n<y\}.
  \]
  The set contains zero and also a positive element, for example
  $t=\min\{1/2,y/2\}$, since $t^n\leq t<y$.
  It is bounded above by $1+y$, because
  $(1+y)^n\geq1+y>y$ and the power map is increasing.
  LUBP therefore gives $x=\sup S>0$.

  Suppose first that $x^n<y$. Choose $h>0$ so small that
  \[
    h<1,\qquad h<\frac{y-x^n}{n(x+1)^{n-1}}.
  \]
  Using \autoref{eq:analysis-i:power-bound} with $a=x$ and $b=x+h$
  gives
  \[
    (x+h)^n-x^n\leq nh(x+h)^{n-1}
      \leq nh(x+1)^{n-1}<y-x^n.
  \]
  Hence $x+h\in S$, contradicting that $x$ is an upper bound.

  Suppose instead that $x^n>y$. Choose $h>0$ with
  \[
    h<x,\qquad h<\frac{x^n-y}{nx^{n-1}}.
  \]
  Now \autoref{eq:analysis-i:power-bound} with $a=x-h$ and $b=x$
  yields $x^n-(x-h)^n\leq nhx^{n-1}<x^n-y$, so $(x-h)^n>y$.
  Every $t\in S$ must then satisfy $t<x-h$, by strict monotonicity.
  Thus $x-h$ is an upper bound of $S$ smaller than its least upper bound,
  again a contradiction. The only remaining possibility is $x^n=y$.
\end{proof}
We may now define $y^{1/n}$ to be this unique nonnegative root.\marginnote{This
proof uses completeness and an algebraic estimate. It does not require
continuity or the intermediate value theorem, which come later in the
course. The two perturbations rule out being below and above $y$.} The inverse
is strictly increasing too. Indeed, if $u<v$ but $u^{1/n}\geq v^{1/n}$,
taking $n$th powers contradicts $u<v$.

The cases $n=2$ and $n=3$ have familiar geometric interpretations.
In $\R^2$, the square $[0,s]\times[0,s]$ has area $s^2$, as shown in
Figure~\ref{fig:analysis-i:root-square}. The theorem supplies its side
length for any prescribed area $A\geq0$. For example, a square of area
$2$ has side length $\sqrt2$.
In $\R^3$, the cube $[0,s]\times[0,s]\times[0,s]$ has volume $s^3$,
as shown in Figure~\ref{fig:analysis-i:root-cube}. A prescribed volume
$V\geq0$ likewise determines its side length, and volume $2$ gives
$s=\sqrt[3]{2}$. These pictures illustrate the root theorem rather than
replace its completeness argument.

\subsection{Exercises on roots and boundaries}
The existence of these roots lets us ask when they can be rational.
For the prime-root exercise, recall Euclid's lemma. If a prime divides
a product of integers, it divides at least one factor; repeated application
gives $p\mid a^k\Longrightarrow p\mid a$.

% !TEX root = ../../hw01.tex
% Homework 01, Exercise 3. Shared problem statement.
\begin{exercise}
  \label{ex:hw01:prime-roots}
  \solutionlink{hw01:prime-roots}
  Let $p=2,3,5,\ldots$ be a prime. Prove that $\sqrt{p}$ is irrational.
  Can this result be extended to $p^{\frac{1}{k}}$ for any integer $k\geq 2$?
\end{exercise}


% !TEX root = ../../hw01.tex
% Homework 01, Exercise 9. Shared problem statement.
\begin{exercise}
  \label{ex:hw01:irrational-roots}
  \solutionlink{hw01:irrational-roots}
  Show that if $a>0$ is irrational, so is the $k$th root of $a$,
  $a^{\frac{1}{k}}$, for any $k\in\mathbb{N}$.
\end{exercise}


The root theorem also lets us identify a Dedekind boundary by a cubic
equation. For the negative side of the cubic, $x<0$ implies $x^3<0$;
on the nonnegative side, use the strict monotonicity just established.
In the next exercise, $A$ is the upper part of the partition; the letters
are reversed from the lower-then-upper convention used in the proof chain.

% !TEX root = ../../hw03.tex
% Source: HW3 Intro to Math Analysis I (1).pdf, page 1, question 2.
% Source correction: part (a) prints B in R, but a partition member is a
% subset of R. The intended set notation is used below.
\begingroup
\renewcommand{\labelenumi}{(\alph{enumi})}
\begin{exercise}
  \label{ex:hw03:cubic-dedekind-cut}
  \solutionlink{hw03:cubic-dedekind-cut}
  Consider $A=\{a\in\R\mid a^3>5-\sqrt{2}\}$.
  \begin{enumerate}
    \item Find $B\subseteq\R$ such that $\{A,B\}$ forms a Dedekind partition.
    \item Find the Dedekind cut for this partition.
    \item Show that $A$ is bounded below and $B$ is bounded above and find
      \[
        \inf A\quad\text{and}\quad\sup B.
      \]
  \end{enumerate}
\end{exercise}
\endgroup


\subsection{A rational exponent must be well defined}
For $x\geq0$ and $m,n\in\N$, define
\[
  x^{m/n}=(x^m)^{1/n}.
\]
This agrees with $(x^{1/n})^m$, since its $n$th power is $x^m$ and it is
nonnegative. We must also check that choosing a different representation
of the fraction gives the same number.

\begin{proposition}
  If $m/n=p/q$ for positive integers $m,n,p,q$, then
  $(x^m)^{1/n}=(x^p)^{1/q}$ for every $x\geq0$.
\end{proposition}
\begin{proof}
  Let $u=(x^m)^{1/n}$ and $v=(x^p)^{1/q}$. Then
  \[
    u^{nq}=x^{mq}=x^{pn}=v^{nq},
  \]
  because $mq=pn$. Both $u$ and $v$ are nonnegative, and the $(nq)$th
  power map is injective there. Therefore $u=v$.
\end{proof}
Thus positive rational powers have a consistent meaning. At zero they
equal zero; for $x>0$, negative rational powers can subsequently be defined
by reciprocals. Real exponents require a further limiting construction.

For the next exercise, $\Q_+$ denotes the positive rational numbers,
and $\R_+=[0,\infty)$ as above. Its second part revisits the equivalent
definitions of rational powers before establishing their algebraic laws.

% !TEX root = ../../hw03.tex
% Source: HW3 Intro to Math Analysis I (1).pdf, page 2, question 6.
\begingroup
\renewcommand{\labelenumi}{(\alph{enumi})}
\begin{exercise}
  \label{ex:hw03:rational-exponent-laws}
  \solutionlink{hw03:rational-exponent-laws}
  For $x\in\R_+$, we defined the function $x^{1/n}$ for $n\in\N$.
  \begin{enumerate}
    \item For $x,y\in\R_+$, prove that $(xy)^{1/n}=x^{1/n}y^{1/n}$.
    \item For $x\in\R_+$, we defined $x^{m/n}=(x^m)^{1/n}$ for
      $m\in\N$. Show that this is the same as $(x^{1/n})^m$.
    \item For $x\in\R_+$ and $r,s\in\Q_+$, prove the laws of exponents
      \[
        (x^r)^s=x^{rs},\qquad x^{r+s}=x^rx^s.
      \]
  \end{enumerate}
\end{exercise}
\endgroup


\subsection{Exercise on odd roots of real numbers}
The preceding root construction was restricted to nonnegative inputs.
The next exercise asks for its extension to the whole real line for odd
powers.

% !TEX root = ../../hw03.tex
% Source: HW3 Intro to Math Analysis I (1).pdf, page 2, question 8.
\begin{exercise}
  \label{ex:hw03:odd-roots}
  \solutionlink{hw03:odd-roots}
  Show that for $n=3,5,7,\ldots$, the function $x^{1/n}$ is a well-defined
  monotone increasing function which is a bijection from $\R$ to itself.
  (Hint: Study $x^n$ as we did in class.)
\end{exercise}


\subsection{Returning to iteration}
We can now finish the square-root argument motivated in
\autoref{sec:analysis-i:convergence-preview}. Let $a>0$, let
$r=\sqrt a>0$ be the root supplied by
\autoref{thm:analysis-i:positive-roots}, and start with $x_0>0$.
Define $x_{n+1}=\tfrac12(x_n+a/x_n)$. Every iterate is positive, and
\[
  x_{n+1}-r=\frac{(x_n-r)^2}{2x_n}\geq0.
\]
Thus $x_n\geq r$ for $n\geq1$, and then
$x_{n+1}-x_n=(a-x_n^2)/(2x_n)\leq0$.
The monotone convergence theorem gives a limit $L\geq r>0$.
The limit laws from \autoref{thm:analysis-i:limit-laws}, applied to
$2x_nx_{n+1}=x_n^2+a$, give $L^2=a$. Root uniqueness therefore gives
$L=r$. If $a$ and $x_0$ are rational, every iterate is rational.
For prime $a$, \autoref{ex:hw01:prime-roots} shows that the limit
is irrational.

The next exercise applies the tangent slopes from
\autoref{sec:analysis-i:convergence-preview} to $F(x)=x-a-1/x$.
The root theorem, limit laws, and monotone convergence are now available
for its fixed-point and convergence arguments. Its continued fraction
denotes the positive solution of $b=a+1/b$.

% !TEX root = ../../hw02.tex
% Homework 02, Exercise 6. Shared problem statement.
\begingroup
\renewcommand{\labelenumi}{(\alph{enumi})}
\begin{exercise}
  \label{ex:hw02:newton-iteration}
  \solutionlink{hw02:newton-iteration}
  In class, we discussed the iteration process
  \[
    x_{n+1}=a+\frac{1}{x_n},\quad n=0,1,2,\ldots,
    \quad a\geq 1,\quad x_0>0,
  \]
  whose limit gives us the continuous fraction
  \[
    b=a+\frac{1}{a+\frac{1}{a+\frac{1}{a+\cdots}}}.
  \]
  That is, $b$ is a fixed point of the function
  \[
    f(x)=a+\frac{1}{x},\quad x>0,
  \]
  which is realized as a root of the function $F(x)=x-f(x)$ $(x>0)$.
  \begin{enumerate}
    \item Show that the iterative algorithm
      \[
        x_{n+1}=\frac{ax_n^2+2x_n}{1+x_n^2},
        \quad n=0,1,2,\ldots,\quad x_0>0,
      \]
      arises from Newton's method that may be used to approach $b$ of our problem.
    \item Assume that $\{x_n\}$ in (a) converges.
      Verify that the limit is indeed the continuous fraction $b$ defined above.
    \item Take $a=\frac{3}{2}$. Prove $0<x_n\leq 2$ for $n=1,2,\ldots$.
    \item Prove that $\{x_n\}_{n=1,2,\ldots}$ increases for any $x_0\ne 2$
      and is constant when $x_0=2$.
    \item Show that $\{x_n\}$ converges to $2$ with any initial state $x_0$
      (thereby establishing the ``global convergence'' of the algorithm).
  \end{enumerate}
\end{exercise}
\endgroup


\subsection{Exercise on Fibonacci ratios}
The final recurrence connects an integer sequence with bounded ratios.
Use induction for its growth and bounds, and the limit laws to identify
a possible limit. The square root appearing in that calculation now
has a rigorous construction. The last part assumes that the ratios
converge; it does not ask for a convergence proof.

% !TEX root = ../../hw03.tex
% Source: HW3 Intro to Math Analysis I (1).pdf, page 1, question 4.
% "Increasing" in part (a) is non-strict, since F_1 = F_2 = 1.
\begingroup
\renewcommand{\labelenumi}{(\alph{enumi})}
\begin{exercise}
  \label{ex:hw03:fibonacci-ratios}
  \solutionlink{hw03:fibonacci-ratios}
  The Fibonacci sequence $\{F_n\}$ is defined by
  \[
    F_0=0,\quad F_1=1,\quad F_n=F_{n-2}+F_{n-1},
    \qquad n=2,3,\ldots.
  \]
  \begin{enumerate}
    \item Show that $\{F_n\}$ is an increasing sequence.
      \emph{Here increasing means nondecreasing.}
    \item Show that $\{F_n\}$ is unbounded.
    \item Show that the Fibonacci ratios
      \[
        r_n=\frac{F_n}{F_{n-1}},\qquad n=2,3,\ldots,
      \]
      is a bounded sequence.
    \item Assume the limit of $\{r_n\}$ exists. Find the limit
      (which happens to be the golden ratio).
  \end{enumerate}
\end{exercise}
\endgroup

